% Transcription — fonds Alexandre Grothendieck, shelfmark 1, batch 8 (pages 141-149).
% Established from the facsimile archives/batches/1/batch-08.pdf (a mirror of
% https://grothendieck.umontpellier.fr/1.pdf), per the skill
% transcribe-grothendieck.
%
% Pass: Fable 5.1 (claude-fable-5-1), 2026-09-10 — first pass, unchecked against
% the pages by a human.
%
% The last batch of the folder: nine pages, all transcribed. \page{} numbers
% the position in the folder, as batches 1 to 7 did: the archivists' pencil
% numbers run two ahead, so page 141 here is pencilled 143 and page 148 is
% pencilled 150.
%
% Three hands and three runs.
%
%   Pages 141-142 — ink, mostly readable: a Proposition giving necessary and
%     sufficient conditions for a sequence of hermitian H_i in 𝓛^p(H) to
%     converge to H_0 (1° ‖H_i‖_p → ‖H_0‖_p, 2° for each λ in the spectrum of
%     H_0 the spectral projections of the H_i on [λ−α, λ+α] eventually have
%     the dimension of that of H_0 and converge), Corollaire 1 (convergence in
%     𝓛^p equivalent to ‖H_i‖_p → ‖H_0‖_p plus convergence in 𝓛^∞),
%     Corollaire 2 (A → |A|^r continuous from 𝓛^p to 𝓛^{p/r}, through A → A*A
%     and H → H^{r/2}), then two Remarques: for 1 ≤ p ≤ 2 is weak convergence
%     plus convergence of norms enough, is 𝓛^p uniformly convex; and is
%     S_p(|H^r − K^r|^{1/r}) increasing in r, which would sharpen
%     Corollaire 1.
%   Pages 143-145 — pencil. Page 143, in ink then pencil, proves the two
%     inequalities ρ_{m+n−1}(vu) ≤ ρ_m(u)ρ_n(v) and ρ_{m+n−1}(u+v) ≤ ρ_m(u) +
%     ρ_n(v) for the singular numbers of operators between Banach spaces,
%     through the subspaces E_{n−1}, F_{m−1} realising ρ_n(u), ρ_m(v); at
%     its foot a question, whether lim sup |λ_i(u)|/ρ_i(u) can be infinite.
%     Pages 144-145, in a faint pencil, restate the convergence results for
%     Banach spaces (spectra σ(u_i) → σ(u), spectral projections), give the
%     continuity of u → |u|^r from L^p(E,F) to L^{p/r}(E) as a Corollaire
%     with the reduction to σ → σ^{r/2} on positive operators, state as
%     Corollaire 2 a criterion for 0 < p ≤ 1, and list questions: uniform
%     convexity of L^p(E,F), continuity of H → H^r for 0 < r < 1 and the
%     inequality ‖|H^r − K^r|^{1/r}‖ ≤ ‖H − K‖, monotonicity of
%     S_{p/r}(H^r − K^r) in r. The connective prose of 144-145 is largely
%     \uncertain{} or \ill{}; the statements carry.
%   Pages 146-147 — pencil, the singular numbers now written r_ν(A): a
%     Théorème (A_i → A in 𝓛^p(H) implies Â_i → Â in Σ_p) reduced to a Lemme,
%     that for every ε there are ρ and i_0 with Σ_{r_ν(A_i) ≤ ρ} r_ν(A_i)^p ≤ ε
%     for i ≥ i_0. The proof introduces the convex function ω(t) equal to e^{pt}
%     for t ≤ log ρ and to its tangent beyond, ω(t) = ρ^p log e(e^t/ρ)^p,
%     reduces to A_i ≥ 0, and splits into a) ρ^p Σ_{r_ν > ρ} log e(r_ν/ρ)^p → 0
%     as ρ → 0, proved by Stieltjes integration against the counting function
%     n(r), ρ^p n(ρ) + pρ^p ∫_ρ^∞ n(r)/r dr; and b), on page 147, from
%     Σ_{r_ν(A_i) ≤ ρ} r_ν(A_i)^p = S_p(A_i) − Σ_{r_ν(A_i) > ρ} r_ν(A_i)^p →
%     S_p(A) − Σ_{r_ν > ρ} r_ν^p. Page 147 also carries a cancelled attempt
%     through Σ ‖A_{i+1} − A_i‖_p < ∞ and A ≤ Σ |A_{i+1} − A_i|.
%   Pages 148-149 — ink, a dense fair hand: the operators P_n(u) =
%     Σ_{k=0}^n α_{n−k}(u) u^k and R(u) = Σ_0^∞ P_k(u) = det(1+u)/(1+u) of a
%     trace-class u, with 1/(1+u) = 1 − u R(u)/det(1+u). Théorème: 1) |P_n(u)| ≤
%     P_n(|u|) and |R(u)| ≤ R(|u|), so P_n(u), R(u) ≥ 0 when u ≥ 0; 2) for u ≥ 0
%     the terms α_{n−k}(u) u^k decrease, whence 0 ≤ P_n(u) ≤ α_n(u)·1 and
%     0 ≤ R(u) ≤ det(1+u); Corollaire (***): ‖R(u)‖_∞ ≤ ‖R(|u|)‖_∞ ≤ det(1+|u|)
%     ≤ e^{‖u‖_1} and ‖P_n(u)‖_∞ ≤ ‖P_n(|u|)‖_∞ ≤ α_n(|u|) ≤ ‖u‖_1^n/n!.
%     The proof rests on a Lemme: for u = Σ λ_i e_i ⊗ e_i', P_n(u) =
%     Σ p_i(u) e_i ⊗ e_i' with p_i(u) = Σ' λ_{i_1}⋯λ_{i_n} over the n-sets not
%     containing i, obtained by computing Σ_{k=0}^n u^{n−k} α_k(u)(−1)^k. A
%     Remarque: the eigenvalues of R(u) are p_i = ∏_{j≠i}(1+λ_j) =
%     det(1+u)/(1+λ_i), whence |p_i| ≤ det(1+|u|) and, since there are λ_i
%     arbitrarily small, ‖R(u)‖_∞ ≥ |det(1+u)| (boxed on page 149), with
%     equality for u ≥ 0. Only the upper third of page 149 is written.
%
% Cancellations: page 141 strikes several words in the statement of the
% Proposition; page 142 carries a pencil note in the lower left margin, struck,
% illegible; page 146 cancels its lines a) and b) with long loops and scribbles
% out one line entirely; page 147 cancels a four-line block with loops. All are
% transcribed as \struck{} with a \note{} where the block cannot be read. The
% vertical pencil notes in the left margins of pages 143 and 146 are recorded
% as \marginal{} with what little can be read.
%
% Notation: 𝓛^p(H) set \mathcal{L}^p(H) on pages 141-142 and 146, as in
% batch 7; the Latin L_0(E), L^p(E,F) of pages 143-145 kept, as in batches 4
% to 6. ρ_n(u) the singular numbers on page 143, r_ν(A) on pages 146-147,
% both as he writes them; S_p(A) = ‖A‖_p^p as in batch 7; |A| = (A*A)^{1/2}.
% « log e(x)^p » on pages 146-147 is his way of writing log(e·x^p), noted once
% on page 146. Sup, Inf, lim sup, det, Tr set as he writes them. « à
% fortiori » and « néc. et suff. » as he spells them. The theorem on page 148
% is numbered by him with stars, (*), (**), (***), kept.
%
% The dating is the inventory's own for the shelfmark. Its group,
% « MATHÉMATIQUES / RECHERCHE / Travaux », carries no date.
\documentclass[11pt,a4paper]{article}
% Le préambule commun à toutes les transcriptions du fonds Grothendieck.
%
% Il tient en une idée : **l'incertitude se compose, elle ne se résout pas.**
% Une transcription qui lisse un mot douteux a détruit la seule chose pour
% laquelle on la faisait — un lecteur doit pouvoir savoir, à chaque endroit,
% ce qui était sur le feuillet et ce qui a été ajouté. Les six macros qui
% suivent sont donc l'appareil critique, et non de la décoration.
%
% Les mêmes commandes sont reconnues par `scripts/render.mjs`, qui produit la
% vue de lecture du site. Une macro ajoutée ici doit l'être aussi là-bas,
% faute de quoi le rendu échouera bruyamment — ce qui est voulu : un rendu
% silencieusement faux est pire qu'un rendu absent.

\ProvidesPackage{grothendieck}[2026/08/08 Transcriptions du fonds Grothendieck]

\usepackage{amsmath,amssymb,amsthm}
\usepackage{mathrsfs}
\usepackage{stmaryrd}
% Les diagrammes commutatifs sont partout dans ce fonds : c'est la langue
% même de la géométrie algébrique de Grothendieck, et une transcription qui
% les rendrait en prose ne transcrirait rien.
\usepackage{tikz-cd}
% Français par défaut — c'est la langue des feuillets — mais l'anglais est
% chargé aussi : la traduction et le résumé sont des documents anglais, et la
% typographie française (espaces avant les deux-points) y serait une faute.
% Ces documents basculent par \selectlanguage{english} en tête de corps.
\usepackage[english,french]{babel}
\usepackage{fontspec}
\usepackage{geometry}
\geometry{a4paper,margin=25mm,marginparwidth=28mm}
\usepackage{marginnote}
\usepackage{xcolor}
% ulem, not soul. `soul`'s \st and \ul are built on a character-by-character
% scan that XeLaTeX's font machinery breaks: the document fails with an
% undefined control sequence rather than degrading. ulem does the same two jobs
% and is Unicode-engine safe. `normalem` stops it hijacking \emph, which the
% transcriptions use for Grothendieck's own emphasis.
\usepackage[normalem]{ulem}
\usepackage{hyperref}
\hypersetup{colorlinks=true,linkcolor=black,urlcolor=black,pdfborder={0 0 0}}

% --- Métadonnées du lot -----------------------------------------------------
% Reprises telles quelles par le script de rendu, qui les affiche en tête de
% la vue de lecture. Elles fixent ce que le fichier prétend couvrir : sans
% elles, une transcription flotte sans point d'attache dans un fonds de seize
% mille feuillets.

\newcommand{\folder}[1]{\gdef\grothFolder{#1}}
\newcommand{\batch}[1]{\gdef\grothBatch{#1}}
\newcommand{\leaves}[2]{\gdef\grothFirst{#1}\gdef\grothLast{#2}}
\newcommand{\foldertitle}[1]{\gdef\grothFoldertitle{#1}}
\gdef\grothFolder{?}\gdef\grothBatch{?}\gdef\grothFirst{?}\gdef\grothLast{?}\gdef\grothFoldertitle{}

% --- Le résumé --------------------------------------------------------------
%
% Il ouvre la lecture modernisée, sous le titre « Résumé », et il est composé
% en petit. Ce n'est pas un ornement : c'est le seul passage du document qui
% ne soit pas des mathématiques, et un lecteur doit pouvoir le reconnaître
% avant de l'avoir lu — puis le sauter s'il connaît déjà le sujet, ou s'y
% arrêter s'il ne le connaît pas. Le filet qui le suit marque la reprise du
% corps.

\newenvironment{resume}
  {\small\setlength{\parskip}{0.45em}}
  {\par\medskip\hrule\medskip}

% --- Mots-clés ---------------------------------------------------------------
%
% En anglais, en clôture du résumé de la lecture modernisée : le vocabulaire
% moderne sous lequel chercher ce que les feuillets construisent. Le manifeste
% du site (`npm run manifest`) les extrait pour étiqueter la cote entière —
% cette ligne est la source unique des tags, il n'y a pas de fichier à côté.
\newcommand{\keywords}[1]{\par\smallskip\noindent
  {\footnotesize\textsc{Keywords} --- \textit{#1}}\par}

% --- L'appareil critique ----------------------------------------------------

% Le numéro de feuillet, en marge. C'est l'ancre qui permet de régler un
% désaccord sur une formule en regardant *un* feuillet plutôt qu'une chemise
% de six cents. Le site s'en sert aussi pour tourner les pages du fac-similé
% quand on fait défiler la transcription.
\newcommand{\leaf}[1]{%
  \marginnote{\footnotesize\color{gray!70}\textsf{#1}}%
  \label{leaf:#1}\ignorespaces}

% Illisible : on ne devine pas. Un mot inventé qui se lit comme les autres est
% le pire résultat possible.
\newcommand{\ill}{\textcolor{red!60!black}{[\,\ldots\,]}}

% Lecture douteuse : le mot est donné, et signalé comme incertain.
\newcommand{\uncertain}[1]{\uline{#1}}

% Ajout de l'éditeur — une lettre manquante, un mot sous-entendu.
\newcommand{\add}[1]{\textcolor{blue!50!black}{[#1]}}

% Biffé par Grothendieck lui-même. Ce qu'il a rayé fait partie du document :
% une rature dit souvent où la pensée a changé de direction.
\newcommand{\struck}[1]{\sout{#1}}

% Note du transcripteur, sur l'état matériel du feuillet.
\newcommand{\note}[1]{\par\smallskip{\small\color{gray!60!black}\textit{[#1]}}\par\smallskip}

% Note marginale de Grothendieck, distincte du corps du texte.
\newcommand{\marginal}[1]{\par\smallskip\noindent\hspace{1em}%
  {\small\color{gray!60!black}#1}\par\smallskip}

% --- Titre ------------------------------------------------------------------

\AtBeginDocument{%
  \begin{center}
    {\large\bfseries Fonds Alexandre Grothendieck --- cote n\textsuperscript{o}~\grothFolder}\\[2pt]
    {\itshape\grothFoldertitle}\\[4pt]
    {\small feuillets \grothFirst--\grothLast\ (lot \grothBatch)}\\[2pt]
    {\footnotesize\color{gray!60!black}Transcription établie d'après le fac-similé numérisé
     par l'Université de Montpellier.\\
     Les lectures incertaines sont soulignées, les illisibles notées [\,\ldots\,],
     les ajouts entre crochets.}
  \end{center}
  \vspace{1em}}

\endinput


\folder{1}
\batch{8}
\pages{141}{149}
\dating{1953}
\watermark{Édition de démonstration}
\foldertitle{Analyse fonctionnelle : notes manuscrites (s.d.), lettre (1953)}

\begin{document}

\note{Dernier lot du dossier : neuf pages, toutes transcrites. Les pages 141 et 142, à l'encre, énoncent une proposition sur la convergence des suites d'hermitiens dans $\mathcal{L}^p(H)$ et ses corollaires. Les pages 143 à 147, au crayon, reprennent la question pour les espaces de Banach, démontrent deux inégalités sur les nombres singuliers, et établissent un lemme sur les petites valeurs singulières d'une suite convergente. Les pages 148 et 149, à l'encre, portent un théorème sur les opérateurs $P_n(u)$ et $R(u) = \det(1+u)/(1+u)$ d'un opérateur à trace.}

\page{141}
\emph{Proposition.} \uncertain{Pour} \uncertain{que} \uncertain{la} \uncertain{suite} $(H_i)$ \uncertain{d'hermitiens} $\in \mathcal{L}^p(H)$ converge vers $H_0 \in \mathcal{L}^p(H)$, il faut et il suffit que

1°) $\|H_i\|_p \to \|H_0\|_p$ (ou seulement $\limsup_{i \to \infty} \|H_i\|_p \leq \|H_0\|_p$)

2°) Pour \struck{tout} $\lambda$ du spectre de $H_0$, (soit $\alpha > 0$ tel qu'il n'existe pas d'autre \struck{$[\lambda - \alpha, \lambda + \alpha]$ ne contienne} valeur spectrale $\mu \neq \lambda$ de $H_0$ telle que $|\lambda - \mu| \leq \alpha$ \struck{\uncertain{l'ensemble} \uncertain{spectral} \uncertain{de} $H$ \ill{}} \struck{et telle} et tel que $\lambda - \alpha$ et $\lambda + \alpha$ n'appartiennent aux spectres d'aucun $H_i$) alors la projection spectrale \struck{$H_i$,} de $H_i$ \struck{relative} dans $[\lambda - \alpha, \lambda + \alpha]$ est de dimension \struck{\ill{}} \struck{égale à} qui finit par être égale à celle de $H_0$, et $h_i \to h_0$ (\uncertain{inutile} \uncertain{de} \uncertain{spécifier} la topologie, \uncertain{cela} \uncertain{résulte} d'une proposition préliminaire sur l'identité de \struck{toutes} \uncertain{ces} topologies) \note{les projections spectrales sont notées $h_i$, $h_0$ en minuscules ; la parenthèse finale est serrée en bas de l'énoncé, sur deux lignes}

\emph{Corollaire 1.} Les deux conditions suivantes sont néc. et suff. pour la convergence de \uncertain{la} \uncertain{suite} $H_i$ vers $H_0$

1) $\|H_i\|_p \to \|H_0\|_p$ (\uncertain{resp.} \ill{})

2) $H_i \to H_0$ dans $\mathcal{L}^\infty(H)$

\emph{Corollaire 2.} $A \to |A|^r$ est \uncertain{application} continue de $\mathcal{L}^p$ dans $\mathcal{L}^{p/r}$.

En effet, elle est composée de $A \to A^{*}A$, continue de $\mathcal{L}^p$ dans $\mathcal{L}^{p/2}$, avec

\page{142}
$H \to H^{r/2}$ de $\mathcal{L}^{+\,p/2}$ dans $\mathcal{L}^{p/r}$ ; or, les conditions de la prop.~1 sont manifestement vérifiées \struck{\ill{}} (condition 1 triviale, condition 2 se voit \struck{ainsi} en \uncertain{se} \uncertain{ramenant} à \struck{\ill{}} \note{en interligne : « la dimension »} \uncertain{finie}, \uncertain{on} \uncertain{est} \uncertain{ramené} \uncertain{au} \uncertain{cas} \uncertain{trivial} ($\sum_{i=1}^{n} \lambda_i\, a_i \otimes a_i$, $\lambda_i \geq 0$ et $a_i$ \uncertain{variables}, $\to \lambda_i^{(0)}$ et $a_i^{(0)}$ ; \ill{} \uncertain{pour} $\sum \lambda_i^{r/2}\, a_i \otimes a_i$ \ill{}) \note{le $\mathcal{L}^{+\,p/2}$ porte un $+$ en exposant, pour le cône des hermitiens positifs}

\emph{Remarque 1} Si $1 \leq p \leq 2$, \uncertain{pour} \uncertain{qu'une} \uncertain{suite} d'hermitiens $H_i$ dans $\mathcal{L}^p$ converge, il faut et il suffit qu'elle converge faiblement, et que $\|K_i\| \to \|K\|$ \note{ainsi, $K$ pour $H$, et la norme sans indice} (\uncertain{cela} \uncertain{est} \uncertain{connu} \uncertain{au} \uncertain{cas} hilbertien) Est-ce vrai si $p < +\infty$ ? \uncertain{Uniforme} convexité de $\mathcal{L}^p$ ?

\emph{Remarque 2} Est-ce que $S_p\bigl(|H^r - K^r|^{1/r}\bigr)$ est fonction croissante de $r$ ($H$ hermitien positif) ? --- Il suffirait de le \uncertain{savoir} pour $r < 1$ \ill{}.
\[
  \bigl\|\,|H^r - K^r|^{1/r}\,\bigr\| \leq \|H - K\|
\]
--- \uncertain{Alors} on aurait une \uncertain{précision} notable \uncertain{dans} le corollaire 1.

\marginal{\ill{}} \note{note au crayon, en biais dans la marge gauche inférieure, biffée de grands traits ; illisible}

\page{143}
\note{le haut de la page est à l'encre, la suite au crayon ; les deux inégalités sont réunies par une accolade}
\begin{gather*}
  \rho_{m+n-1}(vu) \leq \rho_m(u)\, \rho_n(v) \\
  \rho_{m+n-1}(u+v) \leq \rho_m(u) + \rho_n(v)
\end{gather*}
\marginal{$u \in$ \ill{} $E \to F \to G$} \note{dans la marge droite, en regard de la première ligne de la démonstration}

En effet, $\rho_{m+n-1}(vu) = \operatorname*{Inf}_{E_{m+n-2}} \operatorname*{Sup}_{x \perp E_{m+n-2},\ \|x\| \leq 1} \|vux\|$. Or il existe $E_{n-1} \subset E$ tel que $\operatorname*{Sup}_{\|x\| \leq 1,\ x \perp E_{n-1}} \|ux\| = \rho_n(u)$ \struck{et} il existe $F_{m-1} \subset F$ tel que $\operatorname*{Sup}_{\|y\| \leq 1,\ y \perp F_{m-1}} \|vy\| = \rho_m(v)$, soit \uncertain{alors} $H = \bar{u}(F_{m-1})$ \note{« $E_{m+n-2}$ » biffé sous le $H$ ; le $\bar{u}$ se lit aussi $\bar{u}^{-1}$}, --- on
\[
  \operatorname*{Sup}_{x \perp H \wedge E_{n-1},\ \|x\| \leq 1} \|vux\| \leq \rho_n(u)\, \rho_m(v)
\]
or $p - 1 = \operatorname{codim} H \wedge E_{n-1} \geq \operatorname{cod} H + \operatorname{cod} E_{n-1} \geq m - 1 + n - 1 = (m+n-1) - 1$ \note{ainsi ; le « $1$ » après $p$ est surchargé, et les deux signes sont écrits $\geq$, où l'argument demande une majoration} d'où
\[
  \rho_p(vu) \leq \operatorname{Inf} \operatorname*{Sup}_{x \perp E_{p-1},\ \|x\| \leq 1} \|vux\| \leq \rho_n(u)\, \rho_m(v)
\]
à fortiori $\rho_{m+n-1}(vu) \leq \rho_n(u)\, \rho_m(v)$ \note{encadré}.

De même, si $u, v \in L_0(E, F)$, --- on a $p - 1 = \operatorname{codim}(E_{n-1}, F_{m-1}) \geq m - 1 + n - 1 = (m+n-1) - 1$ \note{même remarque sur le sens du signe}, $\rho_p(u+v) \leq \operatorname*{Sup}_{x \perp E_{n-1},\, F_{m-1},\ \|x\| \leq 1} \|(u+v)x\| \leq \rho_n(u) + \rho_m(v)$ \note{le premier membre $\rho_p(u+v) \leq$ est barré d'un trait ou souligné} d'où à fortiori $\rho_{m+n-1}(u+v) \leq \rho_p(u+v) \leq \rho_n(u) + \rho_m(v)$.

\note{sous un trait, en bas de page, au crayon :} \uncertain{Est-ce} \uncertain{qu'on} peut avoir $\limsup \dfrac{|\lambda_i(u)|}{\rho_i(u)} = +\infty$ ; \uncertain{Hypothèse} : \uncertain{invraisemblable} pour $|\lambda_i(u)| \leq \rho_{i/2}(u)$ \uncertain{ou} \uncertain{analogue} ? \note{l'indice du dernier $\rho$ se lit $i/2$}

\marginal{\ill{} \uncertain{parties} \ill{}} \note{trois mots au crayon, soulignés, écrits de bas en haut dans la marge gauche vers le milieu de la page}

\page{144}
\note{page au crayon, pâle ; les énoncés se lisent, la prose entre eux beaucoup moins}

Si $E$ est un Banach, et \uncertain{si} $u_i \to u$ dans $L_0(E)$, alors $\sigma(u_i) \to \sigma(u)$ dans le sens \uncertain{connu} \ill{}. De plus, \uncertain{convergence} des projecteurs \uncertain{spectraux} \ill{}. Dans le cas des op. hermitiens, dans l'esp. de Hilbert, --- \ill{} \uncertain{et} \uncertain{le} \ill{} des conditions nécessaires et suffisantes de convergence.

\uncertain{Conséquence} \uncertain{de} \uncertain{l'hermiticité} \note{en interligne : « $ds \to dl$ », incertain} \uncertain{dans} $L^p(E)$ : \uncertain{deux} \ill{} \uncertain{réciproque} \uncertain{et} \ill{} que (b) \ill{}, (a) $u_i \to u$ \uncertain{spectralement}, i.e. dans $L(E)$ (il suffit \ill{} que l'on ait la convergence spectrale \uncertain{individuelle} en \uncertain{chaque} $\lambda \in \sigma(u)$, $\lambda \neq 0$)

\emph{Corollaire 1} $u \to |u|$ \uncertain{de} $L^p(E, F)$ dans $L^p(E)$ \uncertain{est} continue. Plus \uncertain{généralement}, $u \to |u|^r$ \uncertain{de} $L^p(E,F)$ \uncertain{dans} $L^{p/r}(E)$ \uncertain{est} \uncertain{continue} : \uncertain{c'est} \uncertain{la} \uncertain{composée} \uncertain{de} $u \to u^{*}u$ de $L^p(E,F)$ dans $(L^{p/2}(E))^{+}$, avec $\sigma \to \sigma^{r/2}$ de $(L^{p/2}(E))^{+}$ dans $L^{p/r}(E)$ ; \uncertain{cette} \uncertain{dernière} est continue, car elle est \uncertain{continue} de $L_0(E)^{+}$ dans $L_0(E)$ (\uncertain{ce} \uncertain{qui} \uncertain{se} \uncertain{vérifie} \uncertain{spectralement}, \ill{} \ill{} \ill{} \uncertain{aussi} de $L(E)^{+}$ dans $L(E)$ ? Inégalités précises ?)

\emph{Corollaire 2} \uncertain{Si} $0 < p \leq 1$, pour \uncertain{qu'une} \uncertain{suite} \note{en interligne : « d'hermitiens $u_i \to u$ »} \ill{} converge dans $L^p(E)$, il faut et il suffit que \ill{} $u_i \to u$ \ill{} \uncertain{et} \uncertain{qu'on} \uncertain{ait} \uncertain{la} \uncertain{convergence} \uncertain{faible}.

\page{145}
\note{seul le tiers supérieur de la page est écrit, au crayon pâle}

\uncertain{Question} : \uncertain{Dans} \uncertain{quel} \uncertain{cas} (\uncertain{uniforme} \uncertain{convexité} des \uncertain{esp.} $L^p(E,F)$) \uncertain{vrai} \uncertain{pour} $0 < p < \infty$ ? Pas d'importance.

1) \uncertain{Pour} \uncertain{la} \uncertain{continuité} \uncertain{de} $H \to H^r$, --- \uncertain{se} \uncertain{ramène} \uncertain{au} \uncertain{cas} $0 < r < 1$ (\uncertain{car} \ill{} \ill{} \uncertain{pour} $r$ \uncertain{entier}) \uncertain{Il} \uncertain{semble} \uncertain{probable} \uncertain{que}
\[
  \bigl\|\,|H^r - K^r|^{1/r}\,\bigr\| \leq \|H - K\| \qquad \text{pour } 0 < r < 1 .
\]
\uncertain{Il} \uncertain{revient} \uncertain{au} \uncertain{même} \uncertain{de} \uncertain{prouver} \uncertain{que}
\[
  S_{p/r}(H^r - K^r) = S_p\bigl(|H^r - K^r|^{1/r}\bigr)
\]
est fonction \uncertain{croissante} de $r$, (\uncertain{c'est} \uncertain{le} \ill{} \ill{}, $p$ \uncertain{quelconque}). \uncertain{On} \uncertain{peut} \uncertain{faire} \uncertain{aussi} $p = 1$.

\page{146}
\note{page au crayon ; une note verticale dans la marge gauche, en tête, et une annotation cerclée à côté}

\marginal{\ill{} \uncertain{de} \uncertain{la} \uncertain{continuité} \uncertain{de} $A \to |A|$, \ill{} $|A|$ \ill{} $A^{*}A$ \ill{}} \note{écrite de bas en haut le long de la marge gauche ; à côté, cerclés, trois signes illisibles}

\emph{Th.} Si $A_i \to A$ dans $\mathcal{L}^p(H)$, alors $\hat{A}_i \to \hat{A}$ dans $\Sigma_p$ \note{ainsi : $\hat{A}$ pour la suite des nombres singuliers, $\Sigma_p$ pour l'espace qui les reçoit ; ni l'un ni l'autre n'est défini sur la page}

Cela résulte \uncertain{facilement} \uncertain{du}

\emph{Lemme} Pour tout $\varepsilon > 0$, \uncertain{existe} $\rho > 0$ et $i_0$ tel que \uncertain{pour} $i \geq i_0$ \uncertain{on} \uncertain{ait}
\[
  \sum_{r_\nu(A_i) \leq \rho} \bigl(r_\nu(A_i)\bigr)^p \leq \varepsilon
\]

Pour le \uncertain{démontrer}, soit $\omega(t)$ la fonction \uncertain{convexe} \uncertain{de} $t$ \uncertain{qui} \uncertain{est} \uncertain{égale} \uncertain{à} $e^{pt}$ pour $t \leq \log \rho$, et linéaire pour $t \geq \log \rho$, i.e.
\[
  \omega(t) = p\rho^p\,(t - \log \rho) + \rho^p = \rho^p \log e\Bigl(\frac{e^t}{\rho}\Bigr)^p
\]
\note{« $\log e(x)^p$ » est sa façon d'écrire $\log\bigl(e \cdot x^p\bigr)$, ce que la première expression confirme ; la notation revient jusqu'à la page 147} \ill{} \ill{} (posons $r_\nu = r_\nu(A_i)$) \note{en interligne : « $s_\nu = r_\nu(|A_i|')$ », le dernier signe peu net} \ill{} \uncertain{analogue}) $\sum \omega(\log r_\nu) \leq \sum \omega(\log s_\nu)$ soit
\[
  \sum_{r_\nu \leq \rho} r_\nu^{\,p} + \rho^p \sum_{r_\nu > \rho} \log e\Bigl(\frac{r_\nu}{\rho}\Bigr)^p
  \leq
  \sum_{s_\nu \leq \rho} s_\nu^{\,p} + \rho^p \sum_{s_\nu > \rho} \log e\Bigl(\frac{s_\nu}{\rho}\Bigr)^p
  \qquad \text{(ou } |A_i| \to |A|\text{)}
\]
On peut donc \uncertain{supposer} \uncertain{les} $A_i \geq 0$, (et \ill{} \ill{}) \ill{}, \uncertain{par} \uncertain{extraction} \uncertain{d'une} \ill{} \note{ligne presque entièrement perdue} \ill{} $\varepsilon > 0$, \uncertain{existe} $\rho > 0$ et $i_0$ tel que
\[
  \sum_{r_\nu(A_i) \leq \rho} r_\nu(A_i)^p + \rho^p \sum_{r_\nu(A_i) > \rho} \log e\Bigl(\frac{r_\nu(A_i)}{\rho}\Bigr)^p \leq \varepsilon
\]
\ill{}, \uncertain{posons} $r_\nu = r_\nu(A_i)$ ; il suffit \uncertain{de} \uncertain{montrer} \ill{} :

\struck{a) $\rho^p \sum_{r_\nu > \rho} \log e\bigl(\frac{r_\nu}{\rho}\bigr)^p \to 0$ si $\rho \to 0$}

\struck{b) \uncertain{L'inégalité} \uncertain{du} \uncertain{lemme} \uncertain{pour} \uncertain{les} \uncertain{hermitiens} \uncertain{positifs} $(A_i$ \ill{}} \struck{\ill{}} \note{les deux lignes a) et b) sont barrées de grandes boucles au crayon, et la ligne qui suit b) est entièrement gribouillée ; leurs démonstrations, qui suivent, ne sont pas biffées}

\emph{Preuve de a)}, --- on a en effet
\begin{align*}
  \rho^p \sum_{r_\nu > \rho} \log e\Bigl(\frac{r_\nu}{\rho}\Bigr)^p
  &= \rho^p \int_{r > \rho}^{+\infty} - \log e\Bigl(\frac{r}{\rho}\Bigr)^p\, dn(r) \\
  &= \rho^p \Bigl[- n(r) \log e\Bigl(\frac{r}{\rho}\Bigr)^p\Bigr]_{\rho+0}^{+\infty} + p\rho^p \int_\rho^{+\infty} \frac{n(r)}{r}\, dr \\
  &= \rho^p\, n(\rho) + p\rho^p \int_\rho^{+\infty} \frac{n(r)}{r}\, dr
\end{align*}
\note{le crochet de la deuxième ligne est écrit $[-n(r)\log \ldots]$, l'argument du $\log$ n'étant pas répété ; $n(r)$ est la fonction de comptage des $r_\nu$, comme au lot 6} \uncertain{et} \uncertain{il} \uncertain{est} \uncertain{bien} \uncertain{connu}, \uncertain{puisque} $\sum r_\nu^{\,p} < +\infty$, \uncertain{que} $\rho^p\, n(\rho) \to 0$

\page{147}
\note{seule la moitié supérieure de la page est écrite, au crayon}

et que $\displaystyle\int_0^{+\infty} \frac{n(r)}{r^{1+p}}\, dr < +\infty$ ; \uncertain{partant} $\displaystyle\int_\rho^{+\infty} \frac{n(r)}{r}\, dr < +\infty$, d'où
\[
  p\rho^p \int_\rho^{+\infty} \frac{n(r)}{r}\, dr < p\rho^p \int_{\uncertain{0}}^{+\infty} \frac{n(r)}{r}\, dr \to 0 \quad \text{si } \rho \to 0 .
\]
\note{l'exposant des deux $\rho$ se lit $p$ ou $\alpha$ ; la borne inférieure de la seconde intégrale est peu nette}

\struck{\ill{} \ill{} \ill{} \ill{} \uncertain{que} \ill{} : $\sum \|A_{i+1} - A_i\|_p < +\infty$, \ill{} \ill{} \ill{} \ill{} $\sum S_p(A_{i+1} - A_i) < +\infty$, \uncertain{alors} $A = \sum A_{i+1} - A_i$, \uncertain{d'où} $A \leq \sum |A_{i+1} - A_i|$} \note{bloc de quatre lignes barré de grandes boucles ; le premier $A_{i+1}$ est lui-même surchargé}

\emph{Preuve de b)} \uncertain{On} \uncertain{a}
\[
  \sum_{r_\nu(A_i) \leq \rho} \bigl(r_\nu(A_i)\bigr)^p = S_p(A_i) - \sum_{r_\nu(A_i) > \rho} \bigl(r_\nu(A_i)\bigr)^p
  \to \struck{\leq}\; S_p(A) - \sum_{r_\nu > \rho} r_\nu^{\,p} = \sum_{r_\nu \leq \rho} r_\nu^{\,p}
\]
\note{les signes sous les deux derniers $\sum$ sont surchargés} ; il suffit de prendre $\rho$ tel que $\sum_{r_\nu \leq \rho} r_\nu^{\,p} < \varepsilon$, puis $i$ assez grand.

\page{148}
\note{page à l'encre, d'une écriture serrée ; les formules se lisent, la prose par endroits seulement}

\uncertain{Naturellement}, \uncertain{si} $u$ \uncertain{est} \uncertain{à} \uncertain{trace},
\[
  \frac{1}{1+u} = 1 - u\, \frac{R(u)}{\det(1+u)} \qquad R(u) = \sum_0^\infty P_k(u) = \frac{\det(1+u)}{1+u}
\]
\note{les deux formules sont en tête de page, la première en retrait, la seconde dans la marge droite sur deux lignes}

\emph{Théorème} 1) Soit $u$ \uncertain{un} \uncertain{opérateur} \uncertain{à} \uncertain{trace} \ill{}. \uncertain{Alors}
\[
  \text{(*), (**)}\qquad \struck{\ill{}}\quad |P_n(u)| \leq P_n(|u|),\qquad |R(u)| \leq R(|u|)
\]
\note{les deux inégalités sont soulignées ; entre les numéros et la première, un signe biffé} En \uncertain{particulier}, si $u$ \uncertain{est} $\geq 0$, \uncertain{il} \uncertain{en} \uncertain{est} \uncertain{de} \uncertain{même} \uncertain{de} $P_n(u)$, \uncertain{et} \uncertain{de} $R(u)$

2) Si $u \geq 0$, alors la suite des \uncertain{valeurs} $\alpha_{n-k}(u)\, u^k$ \uncertain{dans} \uncertain{l'expression} \uncertain{de} $P_n(u)$, $P_n(u) = \sum_{k=0}^{n} \alpha_{n-k}(u)\, u^k$, \uncertain{est} \uncertain{une} \uncertain{suite} \emph{décroissante} \uncertain{d'hermitiens}, \ill{} \ill{}. Donc on a \uncertain{en} \uncertain{particulier}, \uncertain{pour} $u \geq 0$ :
\[
  \text{(**)}\qquad \struck{\ill{}}\quad 0 \leq P_n(u) \leq \alpha_n(u) \cdot 1 \quad \text{d'où} \quad \struck{\ill{}}\quad 0 \leq R(u) \leq \det(1+u)
\]
\note{la dernière inégalité est encadrée ; « (**) » repris ici après un bloc biffé ; le « $\cdot 1$ » est l'opérateur identique}

\struck{\ill{}} \emph{Corollaire} \note{un signe en boucle précède le mot} \uncertain{Si} $u$ \uncertain{est} \ill{} \ill{} \ill{} : \uncertain{on} \uncertain{a} \ill{}, \ill{}
\[
  \text{(***)}\qquad
  \begin{cases}
    \|R(u)\|_\infty \leq \|R(|u|)\|_\infty \leq \det(1+|u|) \leq e^{\|u\|_1} \\[4pt]
    \|P_n(u)\|_\infty \leq \|P_n(|u|)\|_\infty \leq \alpha_n(|u|) \leq \dfrac{\|u\|_1^{\,n}}{n!}
  \end{cases}
\]

\emph{Démonstration} \uncertain{Le} \ill{} \ill{} \note{en interligne : « \ill{} \ill{} »}

\emph{Lemme} Si $u = \sum_1^\infty \lambda_i\, e_i \otimes e_i'$ \note{encadré}, alors $P_n(u) = \sum p_i(u)\, e_i \otimes e_i'$ \ill{}, où
\[
  p_i(u) = \sum_{(I)}{}' \lambda_{i_1} \cdots \lambda_{i_n},
\]
$\sum'$ \uncertain{étant} \uncertain{étendue} \uncertain{aux} \ill{} : \struck{\ill{} \ill{}} \note{en interligne, biffé} \uncertain{les} \ill{} \ill{} \ill{} \ill{} \uncertain{ne} \uncertain{contenant} \uncertain{pas} $i$, $I = \{i_1, \ldots, i_n\}$ \ill{} \note{ce paragraphe est accolé d'un trait vertical dans la marge gauche}

Pour le voir, on \uncertain{calcule} \uncertain{successivement} (\uncertain{d'où} \ill{}) \struck{\ill{}} $u^n - \alpha_1(u)\, u^{n-1}$, $u^n - \alpha_1(u)\, u^{n-1} + \alpha_2(u)\, u^{n-2}$, \uncertain{et} \uncertain{l'on} \uncertain{obtient}
\[
  \sum_{k=0}^{n} u^{n-k}\, \alpha_k(u)\, (-1)^k = (-1)^n \sum{}' \lambda_{i_1} \cdots \lambda_{i_n} \ \ill{}
\]
\ill{} \uncertain{le} \ill{} \uncertain{résultat}. \uncertain{par} \uncertain{récurrence} \ill{} d'où \struck{\ill{}} \ill{}. Maintenant 1) \uncertain{devient} \uncertain{évident}. 2) \uncertain{se} \uncertain{vérifie} \ill{} \uncertain{de} \uncertain{façon} \uncertain{immédiate}.

\emph{Remarque} \uncertain{Notons} \uncertain{que} \uncertain{si} \uncertain{les} $p_i$ \uncertain{sont} \uncertain{les} \uncertain{valeurs} \uncertain{propres} \uncertain{de} $R(u)$, \uncertain{on} \uncertain{a} $p_i = \prod_{j \neq i} (1 + \lambda_j)$ \uncertain{d'où} $|p_i| \leq \prod_j (1 + |\lambda_j|) = \det(1 + |u|)$, \uncertain{on} \uncertain{retrouve} \uncertain{une}

\page{149}
\note{seul le tiers supérieur de la page est écrit ; le reste est vierge}

\uncertain{partie} \uncertain{de} (***). \uncertain{Mais} \note{en interligne au-dessus : « \ill{} \uncertain{pas} \ill{} \uncertain{pour}, »} \uncertain{Mais} \uncertain{si} \struck{\ill{}} \ill{} \ill{}, \uncertain{par} \uncertain{suite} $p_i \struck{\ill{}} = \dfrac{\det(1+u)}{1 + \lambda_i}$ \uncertain{$\leq$} d'où
\[
  \|R(u)\|_\infty \geq |\det(1+u)| \qquad (\ill{})
\]
\note{encadré} (car il existe des $\lambda_i$ \uncertain{arbitrairement} \uncertain{petits}, \uncertain{donc} des $p_i$ \uncertain{arbitrairement} \struck{\uncertain{voisins}} \uncertain{voisins} \uncertain{de} $\det(1+u)$ !) \note{« voisins » récrit en interligne au-dessus du mot biffé} \ill{} \struck{\ill{}} \ill{} (\uncertain{pour} l'inégalité (***) a) \ill{} \uncertain{pour} \uncertain{l'évaluation} \ill{}, \uncertain{car} \uncertain{elle} \uncertain{devient} \uncertain{égalité} \uncertain{pour} $u \geq 0$ !

\end{document}
